\documentclass[12pt]{article}
\usepackage[utf8]{inputenc}
\usepackage[T1]{fontenc}
\usepackage[margin=1in]{geometry}
\usepackage{setspace}
\usepackage{xcolor}
\usepackage[colorlinks=true, urlcolor=blue!60!black, linkcolor=blue!60!black]{hyperref}
\setlength{\parindent}{2em}
\setlength{\parskip}{1.25em}
\renewcommand{\baselinestretch}{1.0}
\begin{document}
\begin{center}
{\LARGE \textit{Flast v. Cohen}, 392 U.S. 83 (1968)\par}
\vspace{0.8em}
{\large Case Brief --- Four-Page Edition\par}
\end{center}
\vspace{1em}

\section{Memory Jogger}

Federal taxpayers challenged federal education grants that funded instruction and textbooks in religious schools; an 8--1 Court, narrowing \textit{Frothingham v.\ Mellon}, held that taxpayers have standing when they attack an exercise of Congress's taxing and spending power as exceeding a specific constitutional limit on that power, here the Establishment Clause (\href{https://supreme.justia.com/cases/federal/us/392/83/}{Justia}).

\section{Detailed Case Facts}

Seven federal income taxpayers, led by Florence Flast, sued Secretary of Health, Education, and Welfare Wilbur Cohen to enjoin expenditures under Titles I and II of the Elementary and Secondary Education Act of 1965, which appropriated nearly \$1 billion in grants to state and local educational agencies. Title I funded services for low-income children, including private-school students; Title II funded library resources, textbooks, and instructional materials for public and private schools. Resting standing solely on their taxpayer status, the plaintiffs alleged that federal funds were financing instruction in reading and arithmetic and buying textbooks for religious schools, in violation of the Establishment and Free Exercise Clauses.

\section{Procedural History}

The Government moved to dismiss under \textit{Frothingham v.\ Mellon} (1923), which for 45 years had barred taxpayer standing because a taxpayer's interest in the Treasury is ``comparatively minute and indeterminable.'' A three-judge district court dismissed on \textit{Frothingham}'s authority over Judge Frankel's dissent. The plaintiffs appealed directly, and the Supreme Court reversed on June 10, 1968.

\section{Judicial Votes}

Reversed, \textbf{8--1}.
\begin{itemize}
\item \textbf{Majority:} Chief Justice Warren, joined by Justices Black, Douglas, Brennan, Stewart, White, Fortas, and Marshall.
\item \textbf{Concurring:} Justices Douglas, Stewart, and Fortas, each separately.
\item \textbf{Dissenting:} Justice Harlan.
\end{itemize}

\section{Holding}

Article III does not bar federal taxpayer suits. Taxpayers have standing to challenge a federal statute when they establish (1) a logical link between their taxpayer status and an exercise of congressional power under the Taxing and Spending Clause of Article I, \S\ 8, and (2) a nexus between that status and a specific constitutional limitation on the taxing and spending power, such as the Establishment Clause. \textit{Frothingham} was narrowed, not overruled.

\section{Analysis of Opinions}

\textbf{Majority (Warren).} Commentators had debated whether \textit{Frothingham} was an Article III bar or prudential self-restraint. Rejecting the Government's separation-of-powers argument, Warren held that standing asks only whether the litigant is a ``proper party'' with ``such a personal stake in the outcome of the controversy as to assure that concrete adverseness'' required by \textit{Baker v.\ Carr}; separation-of-powers concerns arise from the issues, not the plaintiff. To screen out generalized grievances, the first nexus excludes an ``incidental expenditure of tax funds in the administration of an essentially regulatory statute,'' and the second requires a specific limitation, not a claim that Congress exceeded its Article I powers generally. The ESEA was a massive spending program, and Madison's \textit{Memorial and Remonstrance}, warning against compelling ``three pence'' for an establishment, showed that the Establishment Clause specifically limits spending. Mrs.\ Frothingham failed the second nexus because she invoked only the Tenth Amendment and due process.

\textbf{Douglas (concurring).} The double-nexus test would prove unsustainable; the Court should overrule \textit{Frothingham} and let taxpayers act as ``private attorneys general.''

\textbf{Stewart and Fortas (concurring).} Stewart stressed that the holding enforces only the Establishment Clause's personal right not to be taxed for religion. Fortas insisted taxpayer status must never become a ``launching pad'' for other attacks.

\textbf{Harlan (dissenting).} Taxpayers who challenge general appropriations assert no personal stake and are ``non-Hohfeldian'' plaintiffs; neither nexus measures one. Courts should not invent public actions without congressional authorization lest the judiciary become a ``Council of Revision.''

\textbf{Later cases.} \textit{Valley Forge} (1982) denied standing for an executive Property Clause transfer; \textit{Hein} (2007) for discretionary executive spending; and \textit{Winn} (2011) for tax credits.

\section{Examples: Future Applications}

\begin{enumerate}
\item \textbf{Same-Side: Sanctuary Repair Grants.} Congress appropriates \$50 million for repair grants to churches and synagogues. Both nexuses are met.
\item \textbf{Same-Side: As-Applied Grant Challenge.} A statute funds adolescent counseling grants, expressly including religious groups, and some grantees teach doctrine. \textit{Bowen v.\ Kendrick} upheld standing for such challenges.
\item \textbf{Opposite-Side: Outside \textit{Flast}'s Window.} A taxpayer attacks military aid under the Declare War Clause (no specific spending limit), or challenges a religious-charity tax credit (\textit{Winn}) or a secretary's prayer rallies from general funds (\textit{Hein}). No standing.
\item \textbf{Opposite-Side: Surplus Building to a Seminary.} The GSA gives a closed federal building to a seminary under a property statute. As in \textit{Valley Forge}, executive Property Clause action fails the first nexus.
\item \textbf{Fence-Sitter: Recovery-Services Worship Component.} An agency funds worship at church-run recovery sites, calling it discretionary administration (\textit{Hein}) while taxpayers say it is part of the statutory grant program (\textit{Bowen}). Standing turns on how the spending is structured.
\end{enumerate}

\section{Critique}

\textbf{Formalism and arbitrariness.} From Harlan to Scalia in \textit{Hein} and Kagan dissenting in \textit{Winn}, critics note the arbitrariness of granting standing when Congress spends \$1 million directly on religious schools, yet denying it when the President spends \$1 million from general appropriations, the Government transfers \$1 million in surplus land, or a State grants \$1 million in tax credits.

\textbf{Tension with \textit{Lujan}.} Warren's claim that standing does not implicate separation of powers has been repudiated by modern doctrine, which in \textit{Lujan v.\ Defenders of Wildlife} treats concrete, particularized injury as an essential separation-of-powers boundary.

\section{Key Quotations}

\begin{quote}
``First, the taxpayer must establish a logical link between that status and the type of legislative enactment attacked. Thus, a taxpayer will be a proper party to allege the unconstitutionality only of exercises of congressional power under the taxing and spending clause of Art. I, \S\ 8, of the Constitution. ... Secondly, the taxpayer must establish a nexus between that status and the precise nature of the constitutional infringement alleged. Under this requirement, the taxpayer must show that the challenged enactment exceeds specific constitutional limitations imposed upon the exercise of the congressional taxing and spending power, and not simply that the enactment is generally beyond the powers delegated to Congress by Art. I, \S\ 8.''
\end{quote}
\noindent --- Chief Justice Warren, Opinion of the Court, 392 U.S. at 102--103.

\begin{quote}
``The taxpayer's allegation in such cases would be that his tax money is being extracted and spent in violation of specific constitutional protections against such abuses of legislative power.''
\end{quote}
\noindent --- Chief Justice Warren, Opinion of the Court, 392 U.S. at 106.

\begin{quote}
``To describe those rights and interests as personal, and to intimate that they are in some unspecified fashion to be differentiated from those of the general public, reduces constitutional standing to a word game played by secret rules.''
\end{quote}
\noindent --- Justice Harlan, dissenting, 392 U.S. at 129.

\end{document}
